\documentclass[10pt,a4paper]{amsart}
\usepackage[margin=19mm]{geometry}
\usepackage{amsmath,amssymb,mathtools}
\usepackage[scr=rsfs,cal=euler]{mathalfa}
\usepackage{xfrac}
\usepackage[unicode,hidelinks,bookmarks=false]{hyperref}
\newcommand{\ie}{i.e.\ }
\newcommand{\cf}{cf.\ }
\newcommand{\resp}{respectively\ }
\newcommand{\Subsection}{\subsection}
\providecommand{\nobreakdash}{\nobreak\hbox{-}}
\setlength{\emergencystretch}{2em}
\multlinegap0pt
\usepackage{amssymb}
\usepackage{enumerate}
\usepackage[shortlabels]{enumitem}
\usepackage{float}
\usepackage{tikz}
\usepackage{mathtools}
\usepackage{braket}
\usepackage[normalem]{ulem}
\usepackage{xcolor}
\newtheorem{lemma}{Lemma}
\title{Positive and negative $3$-energies of graphs}
\author{Zhengbo Chen, Zhouningxin Wang, Xiao-Dong Zhang}
\date{Source: arXiv:2604.15656v1 (CC BY 4.0)}
\begin{document}
\maketitle

\noindent\emph{Source and licence.} Positive and negative $3$-energies of graphs. Version arXiv:2604.15656v1 (CC BY 4.0), distributed by arXiv under the Creative Commons Attribution 4.0 International licence, http://creativecommons.org/licenses/by/4.0/, as recorded in the arXiv licence metadata for this exact version and shown on the abstract page https://arxiv.org/abs/2604.15656v1. Full licence notice: \texttt{licenses/arxiv-2604.15656-CC-BY-4.0.txt}.\par\smallskip

\noindent\textit{Editorial scope.} Counted unit: the article's lemma lem:K\_n-e (source0.tex lines 672--674). The article's complete original proof of it is delivered with it (source0.tex lines 677--704, one \texttt{proof} environment of 28 lines). No mathematics is written, repaired or re-derived: the statement and the proof are the article's own lines, in the article's own notation, and no retained line is substituted or re-flowed.  No further source-local context is needed: every cross-reference of every retained line resolves inside the two delivered spans.  Nothing was omitted from any retained span, so the package carries no omission record.  No retained line contains a citation, so the wrapper carries no bibliography and no bibliography entry.  No retained line contains a figure environment, an image-inclusion command or drawing code, so no image is delivered and the package declares no asset dependency.  Editorial formatting only, in addition: an AMS \texttt{amsart} preamble that restates the article's own declarations and macros used by the retained lines, namely the environments \texttt{lemma} on separate counters, together with the standard packages those lines need.  The retained environments print in this document as Lemma 1 in order of appearance, whereas the article prints the same items with the numbering of its own full document.  The complete native source of the article as distributed in the arXiv e-print for this exact version is bundled unchanged as \texttt{sources/198630/source0.tex}.
\begin{lemma}\label{lem:K_n-e}
Let $n$ be a positive integer with $n\geq 5$. If $G$ is a graph formed from $K_n$ by deleting one edge, then $\mathcal{E}_3^{-}(G)\geq n+1$.
\end{lemma}

\begin{proof}
Assume that $n\geq 5$. Let $V(G)=\{v_1, v_2, \ldots, v_n\}$ and assume that $v_1v_2\not\in E(G)$. We consider a vertex partition $\mathcal{P}$ of $V(G)=\{v_1\}\cup \{v_2\}\cup \{v_3, \ldots, v_n\}$. The adjacency matrix $A(G)$ and the corresponding quotient matrix $M_\mathcal{P}$ of $G/\mathcal{P}$ are as follows:
$$A(G)=
\begin{pmatrix}
0 & 0 & 1 & 1 & \cdots & 1 \\
0 & 0 & 1 & 1 & \cdots & 1 \\
1 & 1 & 0 & 1 & \cdots & 1 \\
1 & 1 & 1 & 0 & \cdots & 1 \\
\vdots & \vdots & \vdots & \vdots & \ddots & \vdots \\
1 & 1 & 1 & 1 & \cdots & 0
\end{pmatrix}
\text{~~and~~} M_\mathcal{P}=
\begin{pmatrix}
0 & 0 & n-2\\
0 & 0 & n-2\\
1 & 1 & n-3
\end{pmatrix}.$$
Observe that for the adjacency matrix $A(G)$, $(0,0,1,-1,0,\ldots,0), (0,0,1,0,-1,\ldots,0), \ldots,$ and $(0,0,1,0,\ldots,0,-1)$ are the eigenvectors of the eigenvalue $-1$ with multiple $n-3$, and $(x,y,z,z,z,\ldots,z)$ are the eigenvectors of other three eigenvalues. So we only consider the quotient matrix $G/\mathcal{P}$.

Note that $\det(\lambda I-M_\mathcal{P})=\lambda^3-(n-3)\lambda^2-(2n-4)\lambda$, which has one zero root and two nonzero roots, denoted by $\lambda_1$ and $\lambda_2$.
It is easy to get that 
$$\lambda_1=\frac{n-3+\sqrt{n^2+2n-7}}{2}>0 \text{~~and~~}
\lambda_2=\frac{n-3-\sqrt{n^2+2n-7}}{2}<0.$$
Thus the negative eigenvalues of $G$ are $-1$ (with multiplicity $n-3$) and $\lambda_2$. 
Then
$$\mathcal{E}_3^-(G)=|-1|^3\times (n-3)+|\lambda_2|^3=n-3+(\frac{\sqrt{n^2+2n-7}-(n-3)}{2})^3 \geq n+1.$$
The last inequality holds because $|\lambda_2|$ is increasing in $n$, and a direct computation shows that $|\lambda_2|^3 > 4$  for all $n \geq 5$.
\end{proof}

\end{document}
